% Redraw of fig-sealed-pillar-pipeline to the tikz-diagram-craft v2 standard.
%
% Reader question: which fences and gates does one work order pass, and which
% checker establishes each resulting property?
% Claim: one work order moves left to right through two fences and then one of
% two release gates; the fences establish enforceability, while the gates make
% noninterference, conservation, and detection meaningful.
% Evidence: the "two fences" / "whole-worker taint and two gates" paragraphs
% and the four-links table, Sec.~sealed-design, in
% website-v2/public/whitepaper/sealed-harbor.tex.
% Must distinguish: fences from gates; the one property established by the
% fences from the three properties that depend on gated release; mechanism
% flow from checker provenance.
% Grammar chosen: a full-width UML component/dataflow rail with paired release
% gates and a compact property strip. The left-to-right direction is the
% runtime order, while direct labels under the rail map properties to checkers.
% Grammar rejected: the former top-to-bottom tree, which spent a full page on
% a narrow stack and made three fan-out arrows look like three separate flows.
% Five-second test: follow C left to right, then name the one fences property
% and any one gates property with its checker.
\begin{figure}[H]
\centering
\pdfigurehue{pdcobalt}%
\begin{tikzpicture}[pd figure,x=1cm,y=1cm]
  \def\flowy{3.05}
  \def\propy{0.05}
  % --- the runtime rail ----------------------------------------------------
  \node[pd artifact,minimum width=2.05cm] (WO) at (0.85,\flowy)
    {work order $C$};
  \node[pd uml component,minimum width=2.45cm] (MON) at (3.55,\flowy)
    {in-guest\\monitor};
  \node[pd uml component,minimum width=2.45cm] (GW) at (6.45,\flowy)
    {outer\\gateway};
  \node[pd uml component,minimum width=2.35cm] (DG) at (9.45,{\flowy+.68})
    {Derek's gate};
  \node[pd uml component,minimum width=2.35cm] (EG) at (9.45,{\flowy-.68})
    {Erin's gate};
  \node[pd focus state,minimum width=2.30cm] (OUT) at (12.55,\flowy)
    {declared release};
  \pdumlicon{MON}\pdumlicon{GW}\pdumlicon{DG}\pdumlicon{EG}
  \draw[pd focus arrow] (WO) -- (MON);
  \draw[pd focus arrow] (MON) -- (GW);
  \draw[pd focus arrow] (GW.east) -- ++(.42,0) |- (DG.west);
  \draw[pd focus arrow] (GW.east) -- ++(.42,0) |- (EG.west);
  \draw[pd focus arrow] (DG.east) -- ++(.42,0) |- (OUT.west);
  \draw[pd focus arrow] (EG.east) -- ++(.42,0) |- (OUT.west);
  \begin{scope}[on background layer]
    \node[pd panel,fit={(MON)(GW)},inner xsep=7pt,inner ysep=11pt] (FENCES) {};
    \node[pd panel,fit={(DG)(EG)},inner xsep=7pt,inner ysep=7pt] (GATES) {};
  \end{scope}
  \node[pd title,anchor=south] at (FENCES.north) {two fences};
  \node[pd title,anchor=south] at (GATES.north) {two gates};

  % --- property-to-checker mapping, directly under the mechanism ----------
  \node[pd kind,anchor=south] at (1.70,0.92) {from the fences};
  \node[pd kind,anchor=south] at (8.80,0.92) {from gated release};
  \node[pd ready state,text width=2.85cm] (PE) at (1.70,\propy)
    {\textbf{Enforceability}\\\texttt{b3\_controllability.py}};
  \node[pd ready state,text width=3.05cm] (PN) at (5.35,\propy)
    {\textbf{Noninterference}\\\texttt{c1\_noninterference.py}};
  \node[pd ready state,text width=2.95cm] (PC) at (9.05,\propy)
    {\textbf{Conservation}\\\texttt{a3\_epsilon\_ledger.py}};
  \node[pd ready state,text width=2.95cm] (PD) at (12.60,\propy)
    {\textbf{Detection}\\\texttt{a4\_canary\_sprt.py}};
\end{tikzpicture}
\caption{A full-width component flow from work order to declared release,
with the four resulting properties keyed directly to their checkers. The two
fences establish enforceability; the two gates make the slot, ledger, and
detector claims testable. [internal; checker provenance named in figure]}
\label{fig:sealed-pillar-pipeline}
\end{figure}
